% G3-5 Validasi walk-forward untuk tuning hyperparameter.
% Sumber: diolah penulis berdasarkan cerqueira2020 (prequential in blocks,
% window bertambah). Istilah lain untuk skema ini sengaja tidak dipakai;
% ikuti penamaan Cerqueira dkk.
\begin{tikzpicture}[x=1mm,y=-1mm,font=\footnotesize,
    sublatih/.style={draw,minimum height=7mm,inner sep=0pt,anchor=west},
    jeda/.style={draw,pattern=north east lines,minimum height=7mm,
      inner sep=0pt,anchor=west},
    valid/.style={draw,fill=black!18,minimum height=7mm,inner sep=1.2pt,
      anchor=west,font=\footnotesize},
    panah/.style={-{Stealth[length=2.2mm]}}]

  % --- penanda: semuanya di dalam data latih ---
  \node[draw,anchor=west,align=center,inner sep=2.5pt,minimum height=7mm,
        minimum width=96mm,fill=black!6] at (16,-12)
    {\textbf{periode pengembangan}};

  \draw[panah] (16,-1) -- (112,-1);
  \node[anchor=south,font=\footnotesize] at (64,-1) {arah waktu};

  % --- 5 fold: jendela latih bertambah, blok validasi bergeser ke kanan ---
  % Data latih membentang x = 16..112 (96mm), dibagi 6 segmen 16mm.
  % Fold i: latih = 16*i - 3 mm, jeda 3mm, validasi 16mm.
  % Ujung kanan fold i = 16 + (16i-3) + 3 + 16 = 32 + 16i; fold 5 -> 112. OK.
  \foreach \i in {1,...,5}{
    \pgfmathsetmacro{\y}{8 + (\i-1)*13}
    \pgfmathsetmacro{\w}{16*\i - 3}
    \pgfmathsetmacro{\xj}{16 + \w}
    \pgfmathsetmacro{\xv}{\xj + 3}
    \node[anchor=east,font=\small] at (14,\y) {\textit{fold} \i};
    \node[sublatih,minimum width=\w mm] at (16,\y) {latih};
    \node[jeda,minimum width=3mm]        at (\xj,\y) {};
    \node[valid,minimum width=16mm]      at (\xv,\y) {validasi};
  }

  % --- label jeda, menunjuk ke blok jeda pada fold 5 ---
  \draw[panah] (84,76) -- (94,66);
  \node[anchor=north east,align=right] at (83,72)
    {jeda $h_b$ bar};

  % --- legenda ---
  \node[sublatih,minimum width=6mm] at (16,90) {};
  \node[anchor=west] at (23,90) {jendela latih};
  \node[jeda,minimum width=6mm] at (56,90) {};
  \node[anchor=west] at (63,90) {jeda};
  \node[valid,minimum width=6mm] at (98,90) {};
  \node[anchor=west] at (105,90) {blok validasi};

\end{tikzpicture}
